\section{近期研究（2026 年 8--11 月）}
\paragraph{从视觉和少量录音估计物体的模态声场}
2026.08~\cite{2608_05145}。这项工作用多视图图像提供物体空间形态的线索，用少量已知接触产生的录音估计模态声学参数。视觉高斯是约束声音合成的空间表示，预测的声音响应还依赖物体的模态参数与敲击位置。

论文评价未见接触位置的撞击声，并展示接触定位等下游任务。少样本指已记录的撞击声较少，方法仍需要对同一物体的视觉观察与物理响应，证据主要对应所测撞击设置。

\paragraph{WildFireGS：在带植被语义的森林高斯上模拟火灾}
2026.08~\cite{2608_11100}。WildFireGS 将航空影像重建的高斯扩展成带植被类别和燃料属性的场景，经过尺度、重力和场景结构预处理后，用粒子模型在高斯表示上模拟热传递与燃料消耗。高斯同时承载可视外观和参与仿真的空间属性。

论文在合成林地控制变量并在真实航空重建场景演示，评价传播趋势与火隔离带和生物量变化。由航拍不完整观测推断的植被结构和燃料参数会影响结果，合成控制实验不等于现场火势预测精度。

\paragraph{LaGSplat：用拉格朗日潜状态驱动动态高斯}
2026.08~\cite{2608_16324}。常规动态高斯以时间索引场景，难以直接响应新激励。LaGSplat 用同一个低维潜向量充当耗散拉格朗日系统的广义坐标和高斯解码器的控制变量。模型从单目视频学习潜状态及其动力学，使原语随物理状态变化而运动。

论文在多个真实系统评价记录运动和片段外外推，并比较实物软体机器人的潜在动力学。视频并不直接监督潜变量，实验展示所测系统的运动与响应，不代表在任意系统中已从单目视频辨识唯一物理参数。

\paragraph{从单目动态视频对齐形变外观与隐式本构模型}
2026.08~\cite{2608_22102}。这项工作结合低秩调整、秩深度几何锚点与本构先验正则，使模拟材料的形变和单目渲染互相约束。前者稳定视觉重建，后者利用现有本构先验限制模拟行为。高斯是可微视觉表示，材料规律由隐式模型学习。

在包含六类材料的合成实验中，作者相对 NeuMA 报告粒子位置 Chamfer 距离平均降低 48\%。在 SpringGaus 的真实数据上另用单目监督验证，采用的评价与具有粒子真值的合成实验不同，不能从合成误差直接推断真实材料属性恢复的精度。

\paragraph{LagrangeGS：从预训练高斯轨迹蒸馏非保守动力学}
2026.08~\cite{2608_22773}。LagrangeGS 先从预训练动态高斯提取每粒子的轨迹，再按对象蒸馏势能、初始速度和非保守力，并用局部刚性对齐避免独立粒子解体。推断可改变学习到的势场以观察反事实运动，或通过负时间步向后积分。

论文在 Dynamic Object 和 Dynamic Indoor Scene 的保留视图与时间窗口评价。前向再反向的平均位置循环误差分别为 1.3\% 和 0.19\%，是数值一致性测量。方法依赖预训练轨迹，解耦粒子假设也限制了碰撞与接触交互。

\paragraph{KnockGS：用一次有控制的响应校准物理高斯}
2026.08~\cite{2608_27365}。KnockGS 固定资产、粒子与 MPM 求解器，只校准材料参数的两个尺度。方法从物体专用响应库选取接近已观测激励的局部样本并进行岭回归，得到的参数冻结写回同一仿真器，再预测未参与拟合的另一组激励。

实验比较最近邻检索、全局回归与固定默认参数，并测试保留的激励方向和幅度。识别建立在已知激励、同一资产和固定仿真合同上，不能直接泛化为任意物体和无控制观测下的材料识别。

\paragraph{PhysMAS：按部件组织材料与运动再做共享域仿真}
2026.09~\cite{2609_07174}。PhysMAS 接收运动描述、四张已标定图像和已对齐的高斯场景。场景模块保持物体及部件身份，材料模块为各部件给出物理配置，再用共享域 MPM 仿真生成动态高斯，避免为每个场景反复做生成梯度优化。

论文比较六类材料场景的语义对齐、感知的物理合理性及参考视频运动一致性。可选本构配置和空间分辨率限制了部件行为与长时互动的覆盖，遮挡下的部件关联也仍可能不稳定。

\paragraph{从一次扫描合成视觉运动策略训练示范}
2026.09~\cite{2609_21112}。该工作从一次 RGB 视频扫描还原背景、各物体的高斯资产，再与按关节配置渲染的机器人模型共同组成场景。物体转移由刚体运动学控制，抓取闭合等接触过渡由学习模型估计，以渲染观测与机器人动作配对形成模仿学习示范。

实验以任务、对象和扫描的组合作为评价单元，每组仿真试验 100 次、实机 10 次。ManiSkill3 的受控对比与真实工作台实验衡量不同条件；平面布局、刚性物体运动和接触模型限制了可覆盖的操作。

\paragraph{GaussPDE：让标量场在静态高斯图上扩散并驱动显示}
2026.09~\cite{2609_27264}。GaussPDE 从静态高斯选择活跃原语，构建稀疏传播图并建立场景相关的导通关系。标量场在图上演化并调节高斯外观，渲染期间高斯几何和图连通性保持固定。

论文用独立合成网格求解验证扩散及跨表面泄漏，并展示三个场景数据集上的完整流程。传播仍依赖重建几何和提升后的视觉线索，结果不能直接推广到任意运动物体或任意物理方程。

\paragraph{PneuTac：用 MPM 和高斯共同模拟软指与触觉图像}
2026.09~\cite{2609_38418}。PneuTac 结合材料点法与高斯渲染模拟气动软指及视觉触觉传感器。使用视觉数据校准后，动作和感知网络可近似部分仿真过程，而经校准的模拟器还用于扩展少量真实示范。

论文以安装 Digit 触觉传感器的气动软指为主要操作实例，并对其他器件做建模验证。实时替代模型与完整 MPM 的适用条件不同，新装置的材料、气动和传感条件仍需重新校准。

\paragraph{用可微的拉格朗日传递连接高斯与 SPH 变形仿真}
2026.10~\cite{2610_04336}。高斯负责从影像重建固体外观，SPH 则用独立离散的材料点满足力学建模需要。该工作在参考构型中建立可微的拉格朗日传递，将 SPH 粒子的位移映射到视觉高斯，并允许影像误差沿该映射影响逆问题求解。

论文分别评价正向运动仿真与从影像反演固体参数。可微连接使参数可以通过影像监督估计，边界条件、各自离散的合理性和高斯包络质量仍是解释误差与参数可辨识性的必要条件。

\paragraph{以接触、运动和应力场连接视频与材料参数推断}
2026.10~\cite{2610_17591}。这项工作从多视图形变视频先推断接触或力、运动与应力等中间场，再预测本构类别和连续材料参数。这把最终参数估计拆解为可检查的物理证据，便于区分追踪准确与材料参数准确。

评价包括参数、稠密场、同模拟器回放与跨协议迁移。默认设置采用三视角 RGB 与物体级划分，屈服应力仅在发生屈服的样本上评价。DeformFieldBench 以合成为主，真实回放缺少已标定力测量和材料真值，不能作为绝对参数精度的证据。

\paragraph{ReStep：以预测终态约束长程粒子轨迹}
2026.10~\cite{2610_18168}。ReStep 先用图式变分模型从初态与物理参数预测合理终态，再用受这两个端点约束的扩散模型补全中间粒子轨迹。与逐步自回归相比，整个预测窗口共享终态约束，高斯渲染用于视觉条件下的物理应用。

实验在 FlexiFall 的落体、倾倒和垂挂场景评价长程误差，并在 Topple 测试视觉物理参数估计。若预测终态错误，它也会作为强条件传到中间轨迹，产生不可信的运动。

\paragraph{用引力约束的动态高斯重建模拟星系}
2026.11~\cite{2611_20089}。这项工作以自发光高斯表示恒星，让亮星沿独立引力轨道运动，并以共享流近似较暗恒星的运动。从多个时刻的带标定视图优化外观与运动时，引力先验约束高斯长时间演化。

实验覆盖八个模拟星系，讨论旋转、角动量、尺寸与渲染质量。实际望远镜对单个星系通常只有单方向和单历元观测，不能提供评估所需的多相机、多时刻数据，因此这组实验说明的是仿真条件下的可行性。

\paragraph{CiGS：从高斯重建表面求解固体内部温度场}
2026.11~\cite{2611_20283}。CiGS 用已配准的可见光和热图重建高斯表面，再提取带热边界观测的点云，直接求解边界积分方程以推断内部温度。深度斜率过滤剔除跨深度断裂的采样点，内部积分对密度常数项单独计算，整个过程不需要表面或体积网格。

作者比较完整覆盖和单视图采样下的精度及耗时，并使用无矩阵的可微求解。由边界观测推断内部温度依赖重建表面和正确的物理边界设定，单视图结果尤其受表面覆盖限制。
